\vspace{-2mm}
\section{Implications}
\label{appendix:Implications}
\vspace{-0.5mm}


\vspace{-1mm}
\paragraph{Implications for Human Simulation} 
Human simulation is a strong tool in various applications of social science~\citep{manning2024automated} and role-playing~\citep{Shanahan2023role,chen2024persona}.
Although plenty of works have adopted LLM agents to simulate human behaviors and interactions~\citep{zhou2023sotopia,gao2023large,xu2024ai}, it is still not clear enough whether LLM agents behave like humans in simulation. Our discovery of behavioral alignment between agent and human trust, which is especially high for GPT-4, provides important empirical evidence to validate the hypothesis that humans' trust behavior, one of the most elemental and critical behaviors in human interaction across society, can effectively be simulated by LLM agents. Our discovery also lays the foundation for human simulations ranging from individual-level interactions to society-level social networks and institutions, where trust plays an essential role. We envision that behavioral alignment will be discovered in more kinds of behaviors beyond trust, and new methods will be developed to enhance behavioral alignment for better human simulation with LLM agents. 





\vspace{-2mm}
\paragraph{Implications for Agent Cooperation}


Many recent works have explored a variety of cooperation mechanisms of LLM agents for tasks such as code generation and mathematical reasoning~\citep{li2023camel,zhang2023exploring,liu2023dynamic}. Nevertheless, the role of trust in LLM agent cooperation remains still unknown. Considering how trust has long been recognized as a vital component for cooperation in Multi-Agent Systems (MAS)~\citep{ramchurn2004trust,burnett2011trust} and across human society~\citep{jones1998experience,kim2022trust,henrich2021origins}, we envision that agent trust can also play an important role in facilitating the effective cooperation of LLM agents. In our study, we have provided ample insights regarding the intrinsic properties of agent trust, which can potentially inspire the design of trust-dependent cooperation mechanisms and enable the collective decision-making and problem-solving of LLM agents.


\vspace{-2mm}
\paragraph{Implications for Human-Agent Collaboration}
Sufficient research has shown the advantage of human-agent collaboration in enabling human-centered collaborative decision-making~\citep{cila2022designing,gao2023towards,mckee2022warmth}. Mutual trust between LLM agents and humans is important for effective human-agent collaboration. Although previous works have begun to study human trust towards LLM agents~\citep{qian2024take}, the trust of LLM agents towards humans, which could recursively impact human trust, is under-explored.
In our study, we shed light on the nuanced preference of agents to trust humans compared with other LLM agents, which can  illustrate the benefits of promoting collaboration between humans and LLM agents. 
In addition, our study has revealed demographic biases of  agent trust towards specific genders and races,  reflecting potential risks involved in collaborating with LLM agents.

\vspace{-2mm}
\paragraph{Implications for the Safety of LLM Agents} 
It has been acknowledged that LLMs achieve human-level performance in a variety of tasks that require high-level cognitive capacities such as memorization, abstraction, comprehension and reasoning, which are believed to be the ``sparks'' of AGI~\citep{bubeck2023sparks}. Meanwhile, there is increasing concern about the potential safety risks of LLM agents when they surpass human capacity~\citep{morris2023levels,feng2024far}. To achieve safety and harmony in a future society where humans and AI agents with superhuman intelligence live together~\citep{tsvetkova2024new}, we need to ensure that AI agents will cooperate, assist and benefit rather than deceive, manipulate or harm humans. Therefore, a better understanding of LLM agent trust behavior can help to maximize their benefit and minimize potential risks  to human society.

























